
One of the books I was reading \cite{Applied} has an introductory section to Macaulay2, so to familiarize myself with the language, I read and recreated given examples. These can be found in the appendix \ref{app:macaulay2} with the prefix "BookExample" or at the GitHub link \cite{Github}. These examples were solid introductions to important concepts such as how to create an ideal or Gröbner bases \ref{lst:book-example-2-1-1}, but were helpful also to familiarize myself with the syntax of the language \ref{lst:book-example-2-2}.The process of creating anonymous, composite functions was unfamiliar and an example demonstrating how to use them can be found in \ref{lst:documentation-conditionals}. However, in order to properly use the language, the next step was to implement an algorithm. 